% Groupes caractéristiques et familles de composés au programme de première (L'essentiel, chapitre 8)
{\renewcommand{\arraystretch}{1.2}\setlength{\tabcolsep}{5pt}\setchemfig{atom sep=1.8em}%
\begin{tabular}{>{\raggedright\arraybackslash\otMedium\fontsize{8.5}{10}\selectfont\color{otNavy}}m{20mm}
  >{\centering\arraybackslash}m{25mm} >{\raggedright\arraybackslash}m{24mm}
  >{\centering\arraybackslash}m{40mm} >{\raggedright\arraybackslash}m{22mm}}
\toprule
\rowcolor{otNavy!6}Famille & Groupe & Nom du groupe & Exemple & Nom \\
\midrule
\rule{0pt}{8mm}alcool & \chemfig{-O-H} & hydroxyle & \chemfig{CH_3-CH_2-OH} & éthanol \\[3mm]
\rule{0pt}{9mm}aldéhyde & \chemfig{-C(=[1]O)-[7]H} & carbonyle, en bout de chaîne & \chemfig{CH_3-CH_2-C(=[1]O)-[7]H} & propanal \\[4mm]
\rule{0pt}{9mm}cétone & \chemfig{-C(=[2]O)-} & carbonyle, dans la chaîne & \chemfig{CH_3-C(=[2]O)-CH_3} & propanone \\[3mm]
\rule{0pt}{9mm}acide carboxylique & \chemfig{-C(=[1]O)-[7]O-[1]H} & carboxyle & \chemfig{CH_3-C(=[1]O)-[7]OH} & acide éthanoïque \\[4mm]
\bottomrule
\end{tabular}}
\par\smallskip\legende{Dans une formule semi-développée, les liaisons avec les atomes d'hydrogène ne sont pas représentées.}
